\documentclass[10pt,a4paper]{amsart}
\usepackage[margin=19mm]{geometry}
\usepackage{amsmath,amssymb,mathtools}
\usepackage[scr=rsfs,cal=euler]{mathalfa}
\usepackage{xfrac}
\usepackage[unicode,hidelinks,bookmarks=false]{hyperref}
\newcommand{\ie}{i.e.\ }
\newcommand{\cf}{cf.\ }
\newcommand{\resp}{respectively\ }
\newcommand{\Subsection}{\subsection}
\providecommand{\nobreakdash}{\nobreak\hbox{-}}
\setlength{\emergencystretch}{2em}
\multlinegap0pt
\usepackage{bm}
\usepackage{bbm}
\usepackage{dsfont}
\usepackage{xspace}
\usepackage{cancel}
\usepackage{mathtools}
\usepackage{amsthm}
\makeatletter
\@ifundefined{mydef}{\newtheorem{mydef}{Mydef}}{}
\@ifundefined{myclaim}{\newtheorem*{myclaim}{Claim}}{}
\@ifundefined{myfact}{\newtheorem{myfact}[mydef]{Fact}}{}
\@ifundefined{mylem}{\newtheorem{mylem}[mydef]{Lemma}}{}
\@ifundefined{mysen}{\newtheorem{mysen}[mydef]{Theorem}}{}
\makeatother
\DeclareMathOperator{\MM}{MM}
\DeclareMathOperator{\cl}{cl}
\newcommand{\dP}{\mathbb{P}}
\newcommand{\uhr}{\upharpoonright}
\title[Cascading Variants of Internal Approachability]{Cascading Variants of Internal Approachability}
\author{Hannes Jakob}
\date{Source: arXiv:2404.18571v2 (CC BY 4.0)}
\begin{document}
\maketitle

\noindent\emph{Source and licence.} Cascading Variants of Internal Approachability. Version arXiv:2404.18571v2 (CC BY 4.0), distributed under the Creative Commons Attribution 4.0 International licence, as recorded in the arXiv licence metadata for this exact version and shown on the abstract page https://arxiv.org/abs/2404.18571v2. Full rights notice: licenses/arxiv-2404.18571-CC-BY-4.0.txt.\par\smallskip

\noindent\textit{Editorial scope.} Counted unit: the Mysen whose source label is \texttt{MMThm} in the article (source0.tex lines 368--370, source label \texttt{MMThm}), delivered together with the article's own complete original proof of it (source0.tex lines 372--429, one \texttt{proof} environment of 58 lines). No mathematics is written, repaired or re-derived: statement and proof are the article's own text line for line. Required context retained uncounted: FactMenas (lines 279--289); FactWoodin (lines 328--330); PSeq (lines 348--354). Every definition, display and label of those lines is retained unchanged. Omitted from this package and not redistributed: 1--278, 290--327, 331--347, 355--367, 371--371, 430--782. Nothing inside a retained excerpt is omitted or altered: every retained body is delivered exactly as the article's lines are written, so no omission receipt of any kind accompanies this package. Citations: the retained excerpts carry no citation command at all, so no bibliography entry of the article is transcribed and no reference is redistributed. Reference adaptation: none was needed and none was made; every reference inside the delivered unit resolves inside it, so no unresolved reference or citation is produced. Printed slips kept literal: no printed word, symbol, hypothesis or number of the counted statement or of its proof was altered. Layout and class: the article's own class is \texttt{amsart} with packages including inputenc, babel, csquotes, amssymb, mathrsfs, hyperref, xcolor, eucal, amsmath, amsthm, imakeidx, breqn, diagbox, verbatim; the wrapper is the lane's pinned \texttt{amsart} editorial wrapper at 10pt on A4, which restates the theorem-like environments the retained text needs and adds the article's own macro definitions that the retained text needs (\texttt{\textbackslash MM}, \texttt{\textbackslash cl}, \texttt{\textbackslash dP}, \texttt{\textbackslash uhr}), with extra wrapper packages (bm, bbm, dsfont, xspace, cancel, mathtools). The delivered numbering is the unit's own. Assets: no retained span contains a figure environment, a graphics-inclusion command, a PDF-page inclusion, a table or a TikZ picture; no image file is copied into this package, the package declares no asset dependency and no image file is redistributed with it. Licence: version arXiv:2404.18571v2 of this article is distributed under the Creative Commons Attribution 4.0 International licence, as recorded in the arXiv licence metadata for this exact version and shown on the abstract page https://arxiv.org/abs/2404.18571v2. A full rights notice is delivered at licenses/arxiv-2404.18571-CC-BY-4.0.txt. Characters: every retained excerpt is pure ASCII, so the delivered engine, XeLaTeX with its default Latin Modern text font, reports no missing character.
	\begin{myfact}[Menas]\label{FactMenas}
		Assume $X\subseteq Y\subseteq Z$ are sets and $\mu$ is a regular cardinal. For $A\subseteq [Y]^{<\mu}$, define $A\uhr X:=\{a\cap X\;|\; a\in A\}$ and $A\uparrow Z:=\{a\in[Z]^{<\mu}\;|\;a\cap Y\in A\}$.
		
		\begin{enumerate}
			\item If $C\subseteq[Y]^{<\mu}$ is a club, there is a function $F\colon[Y]^{<\omega}\to[Y]^{<\mu}$ such that
			$$\cl_F:=\{a\in[Y]^{<\mu}\;|\;\forall y\in[a]^{<\omega}\;F(y)\subseteq a\}\subseteq C$$
			and for any $F\colon[Y]^{<\omega}\to[Y]^{<\mu}$, $\cl_F$ is club in $[Y]^{<\mu}$.
			\item If $C\subseteq[Y]^{<\mu}$ is club in $[Y]^{<\mu}$, then $C\uhr X$ contains a club in $[X]^{<\mu}$ and $C\uparrow Z$ is club in $[Z]^{<\mu}$.
			\item If $S\subseteq[Y]^{<\mu}$ is stationary in $[Y]^{<\mu}$, then $S\uhr X$ is stationary in $[X]^{<\mu}$ and $S\uparrow Z$ is stationary in $[Z]^{<\mu}$.
		\end{enumerate}
	\end{myfact}

	\begin{myfact}[Woodin]\label{FactWoodin}
		Assume $\MM$. Let $\Theta$ be a regular cardinal. Whenever $\dP\in H(\Theta)$ preserves stationary subsets of $\omega_1$ there are stationarily many $N\in[H(\Theta)]^{<\omega_2}$ such that there exists a $\dP$-generic filter over $N$.
	\end{myfact}

	\begin{mylem}\label{PSeq}
		Assume $2^{\omega_2}=\omega_3$. $\dP$ forces the existence of $\subseteq$-increasing and continuous sequences $(a_i)_{i\in\omega_1}$ and $(b_i)_{i\in\omega_1}$ with the following properties:
		\begin{enumerate}
			\item $\bigcup_{i\in\omega_1}a_i=H(\omega_2)^V$ and for all $j<\omega_1$, $(a_i)_{i<j}\in[H(\omega_2)^V]^{<\omega_1}\cap H(\omega_2)^V$.
			\item $\bigcup_{i\in\omega_1}b_i=H(\omega_3)^V$ and for all $i<\omega_1$, $b_i\in[H(\omega_3)^V]^{<\omega_1}\smallsetminus H(\omega_3)^V$.
		\end{enumerate}
	\end{mylem}

	\begin{mysen}\label{MMThm}
		Assume $\MM$ and $2^{\omega_2}=\omega_3$. Then there exist stationarily many $N\in[H(\omega_3)]^{\omega_1}$ such that $N\cap H(\omega_2)$ is internally approachable and $N$ is not internally stationary.
	\end{mysen}

	\begin{proof}
		Fix names $(\dot{a}_i)_{i\in\omega_1}$ and $(\dot{b}_i)_{i\in\omega_1}$ for sequences exemplifying Lemma \ref{PSeq}. Let $C\subseteq[H(\omega_3)]^{\omega_1}$ be club and let $\Theta$ be large enough so that $\dP\in H(\Theta)$. Then $\{N\in[H(\Theta)]^{\omega_1}\;|\;N\cap H(\omega_3)\in C\wedge\omega_1\subseteq N\}$ is club in $[H(\Theta)]^{\omega_1}$ by Fact \ref{FactMenas}. By Fact \ref{FactWoodin} we can find $N\prec H(\Theta)$ with size $\omega_1$ such that the following holds:
		
		\begin{enumerate}
			\item $\dP,(\dot{a}_i)_{i\in\omega_1},(\dot{b}_i)_{i\in\omega_1}\in N$;
			\item there exists a $\dP$-generic filter $G$ over $N$;
			\item $\omega_1\subseteq N$;
			\item $N\cap H(\omega_3)\in C$.
		\end{enumerate}
		
		We will show that $N\cap H(\omega_2)$ is internally approachable and $N\cap H(\omega_3)$ is not internally stationary.
		
		To this end, for $i\in\omega_1$, let $a_i$ consist of all $x\in N\cap H(\omega_2)$ such that for some $p\in G\cap N$, $p\Vdash\check{x}\in\dot{a}_i$. Again for $i\in\omega_1$, let $b_i$ consist of all $x\in N\cap H(\omega_3)$ such that for some $p\in G\cap N$, $p\Vdash\check{x}\in\dot{b}_i$. We will show that the sequences $(a_i)_{i\in\omega_1}$ and $(b_i)_{i\in\omega_1}$ are as desired, starting with $(a_i)_{i\in\omega_1}$.
		
		\begin{myclaim}
			The following holds:
			\begin{enumerate}
				\item The sequence $(a_i)_{i\in\omega_1}$ is increasing and continuous;
				\item for every $i\in\omega_1$, $a_i\in[N\cap H(\omega_2)]^{<\omega_1}$;
				\item $\bigcup_{i\in\omega_1}a_i=N\cap H(\omega_2)$;
				\item for every $j\in\omega_1$, $(a_i)_{i<j}\in N\cap H(\omega_2)$.
			\end{enumerate}
		\end{myclaim}
		
		\begin{proof}
			We first show (1). Clearly, $a_i\subseteq N\cap H(\omega_2)$ by the definition. By Lemma \ref{PSeq}, $\dot{a}_i$ is forced to be countable, so we can let $\dot{f}$ be a name for a bijection from $\omega$ into $\dot{a}_i$. Whenever $x\in a_i$, there is a dense set $D_{x,i}$ of conditions forcing $\dot{f}(\check{n})=\check{x}$ for some $n\in\omega$. In particular, there is $p\in G\cap N$ forcing $\dot{f}(\check{n})=\check{x}$. Since $G$ is a filter, this $n$ is unique, which implies that $a_i$ is countable.
			
			Next up is (2). To see that the sequence is increasing, suppose that $i<j\in\omega_1$ and $x\in a_i$, witnessed by $p\in G\cap N$. Since $\Vdash\dot{a}_i\subseteq\dot{a}_j$ by Lemma \ref{PSeq}, $p$ also witnesses $x\in a_j$. To see that the sequence is continuous, suppose that $j\in\omega_1$ is a limit ordinal and $x\in a_j$, witnessed by $p\in G\cap D$. By Lemma \ref{PSeq}, the set $D_{x,j}$ consisting of all $r\in\dP$ such that $r\perp p$ or $r\leq p$ and $r\Vdash\check{x}\in\dot{a}_i$ for some $i<j$ is dense in $\dP$. By elementarity, $D_{x,j}\in N$ and so there is $r\in D_{x,j}\cap G\cap N$. Since $G$ is a filter, $r$ and $p$ are compatible, so we must have $r\leq p$ and $r\Vdash\check{x}\in\dot{a}_i$, i.e. $x\in a_i$.
			
			Now we show (3). To this end, let $x\in N\cap H(\omega_2)$. Let $D_x$ be the set of all $p\in\dP$ which force $\check{x}\in\dot{a}_i$ for some $i\in\omega_1$. By Lemma \ref{PSeq}, $D_x$ is dense in $\dP$ and it lies in $N$ by elementarity. Since $G$ is $\dP$-generic over $N$, there is $p\in D_x\cap G\cap N$. Thus $p$ witnesses that $x\in a_i$ for some $i\in\omega_1$.
			
			Lastly, we show (4). Let $j<\omega_1$. Let $D_j$ be the set of all $p\in\dP$ which force $(\dot{a}_i)_{i<j}=\check{x}$ for some $x\in H(\omega_2)$. Again be Lemma \ref{PSeq}, $D_j$ is dense in $\dP$. Furthermore, $D_j\in N$ by elementarity. Therefore, there is $p\in D_j\cap G\cap N$. The corresponding $x$ witnessing $p\in D_j$ is in $N$ as well by elementarity.
			
			We will show $x=(a_i)_{i<j}$. Since $p\Vdash\check{x}=(\dot{a}_i)_{i<j}$, $x$ is a sequence of length $j$ consisting of countable sets. We will show that $x(i)=a_i$ for every $i<j$. To this end, let $y\in x(i)$. Then $y\in N$, since $x(i)$ is countable and $p\Vdash\check{y}\in\dot{a}_i$. Thus, $y\in a_i$. On the other hand, let $y\in a_i$. Then there is $q\in G\cap N$ forcing $\check{y}\in\dot{a}_i$. Because $G$ is a filter and generic over $N$, there is $r\leq p,q$ in $G\cap N$. Then $r$ forces $\check{y}\in\check{x}(\check{i})$, which implies $y\in x(i)$.
		\end{proof}
		
		In summary, $(a_i)_{i\in\omega_1}$ witnesses that $N\cap H(\omega_2)$ is internally approachable. We now concern ourselves with $(b_i)_{i\in\omega_1}$.
		
		\begin{myclaim}
			The following holds:
			\begin{enumerate}
				\item $(b_i)_{i\in\omega_1}$ is $\subseteq$-increasing and continuous;
				\item for every $i\in\omega_1$, $b_i\in[N\cap H(\omega_3)]^{<\omega_1}$;
				\item $\bigcup_{i\in\omega_1}b_i=N\cap H(\omega_3)$;
				\item for every $i<\omega_1$, $b_i\notin N\cap H(\omega_3)$.
			\end{enumerate}
		\end{myclaim}
		
		\begin{proof}
			Points (1) through (3) follow using the same arguments as before. So we only need to show that (4) holds.
			
			So we show that $\{b_i\;|\;i\in\omega_1\}$ is disjoint from $N\cap H(\omega_3)$. To this end, let $i\in\omega_1$ and let $x\in N\cap H(\omega_3)$. Let $D_{x,i}$ be the set of all $p\in\dP$ such that for some $y$, $p\Vdash \check{y}\in\check{x}\triangle\dot{b}_i$. Then $D_{x,i}\in N$ by elementarity and $D_{x,i}$ is dense in $\dP$: Let $r\in\dP$ be arbitrary. By Lemma \ref{PSeq}, $\Vdash\dot{b}_i\neq\check{x}$. Thus $r$ forces that there is an element of $V$ which is either in $\dot{b}_i$ and not in $\check{x}$ or vice versa. Therefore we can find $p\leq r$ and $y$ such that $p\Vdash\check{y}\in\check{x}\triangle\dot{b}_i$.
			
			Ergo there is $p\in G\cap D_{x,i}\cap N$. The corresponding $y$ is in $N$ as well by elementarity. Then either $p\Vdash\check{y}\in\check{x}\smallsetminus\dot{b}_i$, in which case $y\in x\smallsetminus b_i$, as no element of $G$ can force $\check{y}\in\dot{b}_i$ or $p\Vdash\check{y}\in\dot{b}_i\smallsetminus\check{x}$, in which case $y\in b_i\smallsetminus x$. In summary, $b_i\neq x$. As $x$ was arbitrary, $b_i\notin N\cap H(\omega_3)$.
		\end{proof}
		
		In particular, $\{b_i\;|\;i\in\omega_1\}$ is club in $[N\cap H(\omega_3)]^{<\omega_1}$ and disjoint from $N\cap H(\omega_3)$, so $N\cap H(\omega_3)$ is not internally stationary.
	\end{proof}

\end{document}
